\section{Experiments}

All experiments are deterministic, CPU-only, and generated by the released scripts. We evaluate target error, interval coverage, width, false confidence, and structural-mismatch sensitivity rather than only in-range $R^2$.

\begin{figure*}[t]
\centering
\begin{minipage}[t]{0.49\textwidth}
\centering
\includegraphics[width=\linewidth]{figures/matched_floor_exponent.pdf}
\textbf{(a) Matched one-power alternatives}
\end{minipage}\hfill
\begin{minipage}[t]{0.49\textwidth}
\centering
\includegraphics[width=\linewidth]{figures/hidden_crossover.pdf}
\textbf{(b) Below-noise spectral crossover}
\end{minipage}
\caption{Excellent pilot agreement does not identify the future. (a) Two ordinary floor-plus-power laws match value and slope. (b) A slower spectral component is initially hidden but dominates excess risk and compute-to-target predictions. Both constructions are exact linear-regression risks.}
\label{fig:limits}
\end{figure*}

\paragraph{Exact lower-bound constructions.}
For the matched pair, nine pilot points in $t\in[-0.5,0.5]$ have Gaussian KL divergence $0.028$, while the target predictions at $t_\star=6.5$ differ by $0.0223$. For the hidden-component construction, a single-power fit has maximum residual only $0.027$ observation standard deviations, yet underestimates the model size needed for excess risks $3\times10^{-5}$ and $10^{-5}$ by factors $13.5$ and $186$, respectively (Figure~\ref{fig:limits}).

\begin{figure*}[t]
\centering
\begin{minipage}[t]{0.49\textwidth}
\centering
\includegraphics[width=\linewidth]{figures/continuous_refinement.pdf}
\textbf{(a) Off-grid certificate refinement}
\end{minipage}\hfill
\begin{minipage}[t]{0.49\textwidth}
\centering
\includegraphics[width=\linewidth]{figures/finite_sample_ols.pdf}
\textbf{(b) Trained finite-sample regression}
\end{minipage}
\caption{The positive theory is numerically checkable. Halving exponent-cell width divides the curvature correction by four. Exact OLS debiasing recovers the population spectral curve, allowing a certificate fit on dimensions $M\le96$ to cover the oracle risk at $M=256$.}
\label{fig:positive}
\end{figure*}

\paragraph{Continuous exponents and trained OLS.}
We generate an off-grid mixture with exponents $(0.337,0.873)$ and extrapolate from ten sizes up to $4096$ to $M_\star=10^6$. Grids of 11, 21, 41, 81, and 161 points all cover the true target. The maximum per-unit curvature correction decreases from $0.537$ to $0.00210$, with the expected factor of four per grid halving; interval width decreases from $0.0238$ to $0.00107$ (Figure~\ref{fig:positive}a).

For the trained bridge, we simulate 5,000 exact OLS risk draws at nine dimensions with $n=8000$. Across dimensions $16$--$256$, the largest relative error between the debiased empirical mean and the oracle spectral risk is $6.1\times10^{-5}$. Using only debiased pilots through $M=96$, the continuous certificate at $M=256$ is $[0.02150,0.02268]$, containing the oracle risk $0.022545$ (Figure~\ref{fig:positive}b).

\paragraph{Noisy synthetic stress test.}
We use nine pilot scales, target $10^6$, known standard error $10^{-3}$, and 30 noise realizations for each of a single power, hidden crossover, and three-component spectrum. Baselines are a constrained floor-plus-power fit with parametric bootstrap, a weak-prior Laplace interval, and a positive two-power point fit. \ScaleCert\ uses the continuous exponent interval $[0.05,1.5]$. We call a miss \emph{false-confident} when its interval width is at most $10\%$ of the observed pilot-loss range.

\begin{table}[t]
\centering
\small
\begin{tabular}{@{}lccc@{}}
\toprule
Curve & Bootstrap & Laplace & \ScaleCert \\
 & Cov./FCR & Cov./FCR & Cov./FCR \\
\midrule
One power & .77/.23 & .97/.03 & 1.00/.00 \\
Hidden crossover & .83/.17 & .87/.13 & 1.00/.00 \\
Three components & .00/1.00 & .00/1.00 & 1.00/.00 \\
\bottomrule
\end{tabular}
\caption{Target coverage and false-confidence rate (FCR) over 30 runs per curve; nominal coverage is $0.90$. The broad certificate is conservative but does not become confidently wrong under hidden structure.}
\label{tab:synthetic}
\end{table}

Table~\ref{tab:synthetic} separates point accuracy from calibrated uncertainty. On the correctly specified one-power family, Laplace intervals are comparatively efficient. Under the three-component alternative, both parametric intervals miss every target and every miss is narrow, whereas \ScaleCert\ covers all targets. Its median widths are $0.0756$, $0.0488$, and $0.0433$ for the three rows, making the cost of robustness explicit.

\paragraph{Public rolling-origin curves.}
We evaluate five neural-machine-translation and four image-classification learning curves released with NSL-PFN \citep{lee2025bayesian}. Three chronological cutoffs per curve produce 27 genuine out-of-range predictions. These files lack independent checkpoint standard errors, so theorem-level Gaussian coverage is inapplicable. Instead we compute the smallest uniform radius making the positive class fit each pilot prefix, then inflate that radius to expose a descriptive coverage--width frontier.

\begin{table}[t]
\centering
\small
\begin{tabular}{@{}crr@{}}
\toprule
Radius & IC cov./width & NMT cov./width \\
\midrule
$1\times$ & .25/.007 & .00/.002 \\
$2\times$ & 1.00/.078 & .60/.097 \\
$3\times$ & 1.00/.133 & .93/.146 \\
$4\times$ & 1.00/.182 & 1.00/.186 \\
\bottomrule
\end{tabular}
\caption{Descriptive public-curve sensitivity. ``Radius'' multiplies the minimum in-range structural residual; widths are medians. The result is a model-mismatch frontier, not a distribution-free coverage claim.}
\label{tab:public}
\end{table}

The one-power bootstrap covers $58\%$ of image-classification targets and $27\%$ of NMT targets. Table~\ref{tab:public} shows why a single optimistic certificate is insufficient: the smallest in-range-compatible radius extrapolates poorly, whereas explicit mismatch inflation raises coverage while widening intervals. This is the intended behavior of partial identification---assumptions that cannot be validated from the prefix remain visible in the output.
